\begin{textsample}{POS Dim 1 – es – Score 37.00 – es/es\_ef\_anos\_finais\_linguagens\_ingles\_2.txt}  \label{ex:f1_pos_200}
Na Pedagogia dos Multiletramentos são enfatizados processos de aprendizagem que levam os estudantes a : - Experienciar - O conhecido e o novo, envolvendo a reflexão sobre as práticas letradas já conhecidas dos estudantes, bem como as práticas de \textbf{letramento} que ainda não lhes são familiares, mas que fazem parte de seus universos e que possibilitem novas formas de significação e de produção de sentidos. - Conceitualizar - Por nomeação e com \textbf{teoria}, referindo -se \textbf{tanto} a processos de categorização de elementos \textbf{teóricos}, nomeando -os e conceituando -os, como também a processos mais autônomos e de generalização, tais como criar \textbf{quadros} \textbf{teóricos}, interpretativos e explicativos de determinados fenômenos com a ajuda dos conceitos e nomes dados, tornando os estudantes conceituadores, criadores de conceitos. - Analisar - Processos de interpretação e de \textbf{compreensão} de enunciados, envolvendo \textbf{tanto} elementos de \textbf{análise} \textbf{funcional} da linguagem, como causa e efeito, inferências e deduções, como também elementos de \textbf{análise} \textbf{crítica}, uma vez que \textbf{demandam} do estudante um \textbf{posicionamento} \textbf{crítico} diante dos enunciados aos quais estão sendo expostos, como, por exemplo, refletir sobre a intencionalidade de enunciados e apreciações de valor. - Aplicar - Os conhecimentos construídos ao mundo real de modo criativo e inovador, evidenciando as transformações promovidas no estudante durante as etapas dos processos de aprendizagem.
Para dar conta dessa proposta, de modo a desenvolver as \textbf{Competências} e Habilidades específicas propostas no ensino de Língua \textbf{Inglesa} - Ensino Fundamental, os Objetos de Conhecimento devem ser trabalhados por meio de situações comunicativas contextualizadas com o \textbf{foco} na língua em uso, utilizando -se de práticas de linguagem orientadas por Campos Temáticos que foram organizados nos Eixos de Práticas de Linguagem : oralidade, leitura, escrita, conhecimentos linguísticos e dimensão \textbf{intercultural}, previamente descritos na BNCC.
Assim, o percurso formativo \textbf{sugerido} neste documento está \textbf{detalhado} por ano nos \textbf{quadros} de \textbf{Sistematização} das Aprendizagens Essenciais, com a seguinte estrutura : CAMPO TEMÁTICO, OBJETOS DE CONHECIMENTO, HABILIDADES, \textbf{COMPETÊNCIAS} ESPECÍFICAS E TEMAS INTEGRADORES, a saber : - CAMPO TEMÁTICO : trata -se das práticas de linguagem organizadas nos seguintes Eixos : - Eixo Oralidade - Práticas de \textbf{compreensão} e produção oral de língua \textbf{inglesa}, em diferentes contextos \textbf{discursivos}, presenciais ou simulados, com repertório de falas diversas, incluída a fala do professor. - Eixo Leitura - Práticas de leitura de textos diversos em língua \textbf{inglesa} ( verbais, verbo-visuais, multimodais ) presentes em diferentes suportes e esferas de circulação.
Tais práticas envolvem articulação com os conhecimentos prévios dos \textbf{alunos} em língua materna e/ou outras línguas.
Compreendem as estratégias de leitura, práticas de leitura e construção de repertório lexical, atitudes e disposições do leitor, práticas de leitura e pesquisa, fruição e avaliação de textos lidos. - Eixo Escrita - Práticas de produção de textos em língua \textbf{inglesa}, relacionados ao cotidiano dos \textbf{alunos}, presentes em diferentes suportes e esferas de circulação.
Tais práticas envolvem a escrita mediada pelo professor ou colegas e articulada com os conhecimentos prévios dos \textbf{alunos} em língua materna e/ou outras línguas, especialmente a língua \textbf{inglesa}.
Compreendem estratégias de pré-escrita e escrita, práticas de escrita e produção de escrita. - Eixo Conhecimentos Linguísticos - Práticas de \textbf{análise} linguística para a reflexão sobre o funcionamento da língua \textbf{inglesa}, com base nos usos de linguagem trabalhados nos eixos Oralidade, Leitura, Escrita e Dimensão \textbf{Intercultural}.
Compreendem o estudo do léxico e gramática. - Eixo Dimensão \textbf{Intercultural} - Reflexão sobre aspectos relativos à interação entre culturas ( dos estudantes e aquelas relacionadas a demais falantes de língua \textbf{inglesa} ), de modo a favorecer o convívio, o respeito, a superação de conflitos e a valorização da diversidade entre os povos.
Compreendem os seguintes elementos de reflexão : a língua \textbf{inglesa} no mundo, no cotidiano da sociedade brasileira e na comunidade, comunicação \textbf{intercultural} e manifestações culturais.
Obs. : Todos os Eixos devem ser orientados pelos temas \textbf{sugeridos} para cada ano, indicados nos \textbf{quadros} de \textbf{sistematização} das aprendizagens essenciais, conforme \textbf{segue} : - 6º Ano : “ Escola e seu entorno ” ; “ Família e comunidade ” e “ Diversidade linguística ” ; - 7º Ano : “ Que histórias temos para contar ” ; “ Meu lugar no mundo ” e “ Indústria cultural ” ; - 8º Ano : “ Direitos humanos hoje e amanhã ” ; “ Mundo Digital e o futuro ” ; “ Diversidade cultural ” ; - 9º Ano : “ Sociedade e consumo ” ; “ Internet : suas linguagens e armadilhas ” ; “ Mundo \textbf{globalizado} : identidades ”.
OBJETOS DE CONHECIMENTO : são entendidos como os conteúdos, conceitos e processos que estão relacionados aos Campos Temáticos.
HABILIDADES : expressam as aprendizagens essenciais que devem ser asseguradas aos \textbf{alunos} nos diferentes contextos escolares e devem ser trabalhadas de modo a promover sua \textbf{progressão} ao longo do ano letivo, ou dos anos de estudo, considerando a complexidade dos processos cognitivos envolvidos em cada etapa do processo de aprendizagem. \textbf{COMPETÊNCIAS} ESPECÍFICAS : evidenciam como as Dez \textbf{Competências} Gerais apresentadas na BNCC se expressam nas áreas de conhecimento e nos componentes curriculares.
TEMAS INTEGRADORES : compreendem aspectos que ultrapassam a dimensão cognitiva, de modo a promover a \textbf{compreensão} dos processos sociais e a mudança de comportamentos que comprometem a convivência democrática.
É facultado ao professor integrar aos seus planos de \textbf{aula} outros temas que \textbf{julgue} ser convenientes, considerando as especificidades de seus contextos.
Outro aspecto importante do Currículo do Espírito Santo que deve ser evidenciado no planejamento das \textbf{aulas} de Língua \textbf{Inglesa} é a intencionalidade pedagógica nos processos de desenvolvimento das \textbf{competências} \textbf{socioemocionais}, o que potencializa a construção do repertório linguístico \textbf{discursivo} do estudante, promovendo sua autonomia, a elevação de sua autoestima e autoconfiança, na medida em que oferece possibilidades de aprender a conhecer, aprender a ser, aprender a fazer e aprender a conviver por meio de \textbf{abordagens} que o levem a desenvolver o pensamento \textbf{crítico}, autônomo, participativo, solidário e atuante em todo o processo educacional.
É \textbf{oportuno} assinalar acerca do papel \textbf{primordial} que o professor exerce durante a implementação de uma proposta curricular.
Isso significa \textbf{dizer} que é pela mediação desse importante agente que o currículo acontece na sala de \textbf{aula}.
Sendo o currículo uma prática social, é por meio da relação professor-estudante-conhecimento que as práticas curriculares se concretizam.
Portanto, evidencia -se outro grande desafio para o professor de língua \textbf{inglesa} do Espírito Santo : refletir sobre a sua prática, sobre suas estratégias de ensino, analisando se as atividades propostas aos estudantes condizem com os níveis cognitivos \textbf{demandados} nas “ habilidades ” trabalhadas, se as atividades contemplam as “ etapas dos processos de aprendizagem ” e demais elementos pedagógicos aqui \textbf{sugeridos}, e constantemente refletir se essas práticas curriculares promoveram alguma transformação nos estudantes, \textbf{tanto} no que \textbf{diz} respeito aos aspectos cognitivos quanto aos aspectos \textbf{socioemocionais}.
Os desafios para a concretização dessas transformações são \textbf{inúmeros}, mas aproveitamo -nos da deixa do compositor, e “ last but not least ”, na cabeça \textbf{aberta}, global, conectada, multicultural, ativa, criativa e \textbf{reflexiva} do professor de língua \textbf{inglesa}, “ toda hora rola um insight ”.
Portanto, é \textbf{imprescindível} o entendimento de que o seu “ savoir-fare ” é essencial na profissão que exerce e que o seu “ approach ” é determinante para as transformações no processo de \textbf{ensino-aprendizagem} de Língua \textbf{Inglesa} e na educação aqui almejadas.

% matched lemmas: aberto, abordagem, aluno, análise, aula, competência, compreensão, crítico, demandar, detalhar, discursivo, dizer, ensino-aprendizagem, foco, funcional, globalizado, imprescindível, inglês, intercultural, inúmero, julgar, letramento, oportuno, posicionamento, primordial, progressão, quadro, reflexivo, segar, sistematização, socioemocional, sugerir, tanto, teoria, teórico
\end{textsample}
